% BHU PYQ -- Professional Exam Template
\documentclass[11pt,a4paper]{article}

\usepackage[a4paper,top=2.5cm,bottom=2.5cm,left=2.8cm,right=2.8cm,headheight=14pt]{geometry}
\usepackage{amsmath,amssymb}
\usepackage[T1]{fontenc}
\usepackage[utf8]{inputenc}
\usepackage{lmodern}
\usepackage{microtype}
\usepackage{fancyhdr}
\usepackage{enumitem}
\usepackage{listings}
\usepackage{hyperref}
\usepackage{lastpage}
\usepackage{parskip}

\hypersetup{colorlinks=true,urlcolor=black,linkcolor=black,
  pdftitle={BVCA-301: Analog Digital Communication -- BHU PYQ}}

\pagestyle{fancy}
\fancyhf{}
\renewcommand{\headrulewidth}{0.4pt}
\renewcommand{\footrulewidth}{0.4pt}
\fancyhead[L]{\small\textit{Banaras Hindu University}}
\fancyhead[R]{\small\textit{BVCA-301: Analog Digital Communication}}
\fancyfoot[C]{\small Page \thepage\ of \pageref{LastPage}}

\lstset{
  basicstyle=\small\ttfamily,
  frame=single,
  breaklines=true,
  columns=flexible,
  keepspaces=true
}

\newlist{parts}{enumerate}{2}
\setlist[parts,1]{label=(\alph*),leftmargin=*,itemsep=4pt,topsep=3pt}
\setlist[parts,2]{label=(\roman*),leftmargin=*,itemsep=2pt,topsep=2pt}
\newcommand{\pts}[1]{\hfill\textit{[#1]}}

\begin{document}
\begin{center}
  {\large\bfseries BANARAS HINDU UNIVERSITY}\\[2pt]
  {\normalsize B.Voc. Semester III Examination, 2022-23}\\[6pt]
  \rule{0.8\linewidth}{0.5pt}\\[6pt]
  {\large\bfseries Paper: BVCA-301 --- Analog Digital Communication}\\[4pt]
  {\normalsize Time: Three Hours \hfill Maximum Marks: 70}
\end{center}

\rule{\linewidth}{0.4pt}

\smallskip
\noindent\textit{Instructions: Attempt all questions. Write your answers in your own words as far as practicable.}

\bigskip

\bigskip\noindent{\large\bfseries SECTION --- A}
\rule{\linewidth}{0.4pt}\smallskip

Long-answer questions to be answered in about 750 words each. Attempt any two: \hfill \textbf{15 $\times$ 2 = 30 Marks}

\medskip
\noindent\textbf{Question 1.} Describe Analog and Digital Communication systems. Discuss the fundamentals of Digital Modulation.
\centerline{\textbf{OR}}
Describe 8-QAM and 16-QAM with suitable examples.

\medskip
\noindent\textbf{Question 2.} Describe the detection of signals in AWGN.
\centerline{\textbf{OR}}
Describe Error Correcting Codes with reference to Convolutional Codes.

\bigskip\noindent{\large\bfseries SECTION --- B}
\rule{\linewidth}{0.4pt}\smallskip

Short-answer questions to be answered in about 200 words each. Attempt any six: \hfill \textbf{5 $\times$ 6 = 30 Marks}
\begin{parts}
  \item Describe the classification of Signals.
  \item Discuss PAM with examples.
  \item Describe Linear Block Codes with examples.
  \item What is Inter-Symbol Interference (ISI)?
  \item Describe different wireless standards.
  \item Discuss IEEE 802.11a and 802.11n.
  \item Describe the classification of systems.
  \item Discuss Hamming Codes.
\end{parts}

\bigskip\noindent{\large\bfseries SECTION --- C}
\rule{\linewidth}{0.4pt}\smallskip

Answer all parts in a word or sentence each: \hfill \textbf{10 $\times$ 1 = 10 Marks}
\begin{parts}
  \item Full form of AWGN.
  \item Define Parity.
  \item Is the set of codewords $\{0000, 0101, 1010, 1111\}$ cyclic?
  \item Define Network Communication.
  \item What is the function of an Input Transducer?
  \item Define the BIBO stability criterion.
  \item Define Phase Modulation.
  \item Define Pulse Time Modulation.
  \item Define a Phasor Diagram.
  \item Define a Constellation Diagram.
\end{parts}

\vfill
\rule{\linewidth}{0.4pt}
\begin{center}\small
\href{https://bhu-exam.vercel.app/}{Click here to visit our website}\\[4pt]\href{https://me-aryan.vercel.app/}{Click here to visit our contributor}\\[4pt]\href{https://quashan-me.vercel.app/}{Click here to visit our contributor}
\end{center}
\end{document}
